\documentclass[a4paper,11pt,dvipdfmx]{jsarticle}

\usepackage[top=18mm,bottom=18mm,left=22mm,right=22mm]{geometry}
\usepackage{amssymb}
\usepackage{listings}
\usepackage{plistings}
\usepackage{titlesec}
\usepackage{enumitem}
\usepackage[dvipdfmx,hidelinks]{hyperref}

\titlespacing*{\section}{0pt}{0.5em}{0.1em}
\setlength{\parskip}{0pt}
\setlist[itemize]{itemsep=1pt,topsep=2pt,parsep=0pt,partopsep=0pt}

\lstset{
  basicstyle=\ttfamily\small,
  frame=single,
  breaklines=true,
  xleftmargin=1em,
  aboveskip=6pt, belowskip=6pt,
  columns=flexible,
}

\newcommand{\cboxon}{$\boxtimes$}
\newcommand{\cboxoff}{$\square$}

\begin{document}

{\LARGE\bfseries 生成AI利用申告書}

\begin{itemize}
  \item 氏名：塩澤 拓海
  \item 学籍番号：25G1065
  \item プログラムファイル名：\texttt{HCK02\_01.ino}
  \item 使用したAIツール：Claude Code (Opus 4.7)
\end{itemize}

\section*{1. 何に使ったか}
例題2 のプログラムをそのまま動かすと、シリアルプロッタの縦軸が自動で動いて
しまい、波形が見にくかった。出力のしかたをどう直せばよいか、生成AIに相談した。

\section*{2. AIへの指示（プロンプト）}
\begin{lstlisting}
例題2 のコードをシリアルプロッタで見ると、縦軸の範囲が勝手に変わって
波形が読めません。マイクの値を読み取って、波形が見やすく出るように
出力のしかたを直したいです。どういう出力にすればよいですか?
\end{lstlisting}

\section*{3. AIの出力の使い方}
\begin{itemize}
  \item \cboxoff そのまま使った
  \item \cboxon 一部修正して使った
  \item \cboxoff 参考にしただけ
\end{itemize}
$\rightarrow$ 上で選んだ理由：

AIから出てきた「ラベル付きの文字列を消して数値だけをカンマ区切りで出す」
「縦軸を固定するために毎回 0 と 1023 のガイド線も一緒に出す」という
2点を採用した。AIのサンプルにあった電圧変換などの処理は今回不要だったので
削った。

\section*{4. 正しいと確認した方法}
\begin{itemize}
  \item シリアルモニタを 921600\,bps で開き、
        \texttt{<生値>,0,1023} の形式で1行ずつ流れていることを確認した
  \item シリアルプロッタで上下に \texttt{0} と \texttt{1023} のガイド線が出て、
        中央の線がマイクに入る音に応じて動くことを観察した
\end{itemize}

\section*{5. 問題や間違い}
\begin{itemize}
  \item \cboxon あった
  \item \cboxoff なかった
\end{itemize}

AI が最初に出してきたコードは、\texttt{Serial.print} で各値の前に
\texttt{"value:"} のような文字列ラベルを付ける形だったが、送信バイト数が
増えてプロッタの更新が重くなりそうだったので、整数だけをカンマで区切る
形式に直した。

\section*{6. 自分で工夫した点}
\begin{itemize}
  \item 例題2 にあった電圧換算や中心化処理は今回不要だったので削除し、
        必要最小限の処理だけにした
  \item 毎ループで \texttt{0} と \texttt{1023} を一緒に送り、シリアルプロッタの
        縦軸が動かないようにした
  \item サンプリングは \texttt{delayMicroseconds(100)} で素直な書き方にした
\end{itemize}

\section*{参考文献}
\begin{itemize}
  \item Arduino 日本語リファレンス（武蔵野電波） \\
        \url{https://www.musashinodenpa.com/arduino/ref/}
\end{itemize}

\end{document}
